\section{Introduction}
\label{sec:intro}

Traffic heterogeneity has become a central challenge for backbone networks. The same infrastructure must simultaneously carry interactive video, remote control, distributed AI training, and bulk backup, and these services differ sharply in their preferences for bandwidth, latency, and packet loss across multiple QoS dimensions. The result is a coexistence of mouse flows (low bandwidth, strict QoS requirements) and elephant flows (high bandwidth, best-effort tolerant)~\cite{CISCO2023_GlobalTraffic}. Such heterogeneous demand leaves traditional undifferentiated routing strategies, whether shortest-hop-count OSPF or equal-cost multipath ECMP, unable to satisfy resource utilization and QoS guarantees at the same time.

The fundamental difficulty lies in an online causal coupling: flows arrive sequentially, and the future is unseen. When selecting a route for the $t$-th flow, the agent cannot foresee the subsequent demands $\mathbf{f}_{t+1}, \dots, \mathbf{f}_W$~\cite{Wang1996_QoSRouting}. Were the entire flow sequence available offline, route selection could be cast as an integer program with a globally optimal solution; the online constraint, however, forces decisions under incomplete information, so that earlier and later flows compete for shared link capacity. Greedy heuristics such as CSPF, though locally optimal at each step, frequently place early-arriving mouse flows on resource-rich optimal paths and thereby deplete the link capacity later needed by elephant flows, a classic dilemma of myopic strategies~\cite{Fortz2000_TrafficEngineering}. An optimal policy must instead reserve capacity non-myopically in expectation over the demand distribution, trade off across flows, and attribute the cumulative utility realized at the end of the sequence back to earlier steps. This structure, sequential decision-making under an unseen future with long-range credit assignment, is naturally suited to reinforcement learning (RL).

Combining deep reinforcement learning (DRL) with graph neural networks (GNNs) has opened new possibilities for online routing: DRL supplies non-myopic sequential decision-making~\cite{Sun2022_ScalableRouting_TON,He2024_RoutingOptimization_TMC}, and GNNs provide explicit encoding of network topology~\cite{Rusek2020_JSAC_RouteNet,FerriolGalmes_2023_RouteNetFermi,Munikoti2024_DRL_GNN_TNNLS}. Existing methods nonetheless suffer from two fundamental limitations. First, DRL routing policies typically optimize a single global objective for all flows and cannot provide differentiated service under mixed traffic. Second, GNN-based methods learn one unified topology embedding for all flows and cannot tailor that embedding to per-flow QoS characteristics, the ``same graph, same embedding'' problem.

Breaking this homogeneity constraint requires a network encoder with flow-conditioning capability: it must adjust its topological perception dynamically according to each flow's individual QoS characteristics. Conditioned neural networks, such as FiLM for general networks~\cite{Perez2018_FiLM} and its graph extension GNN-FiLM~\cite{Brockschmidt2020_GNNFiLM}, supply the necessary mechanisms, yet they have not been applied to per-flow QoS-aware routing. In the online routing scenario, what truly needs conditioning is not intra-graph node attributes but per-flow QoS intent: the same topology, when facing a mouse flow versus an elephant flow, should undergo a message-passing process reshaped into two entirely different topological perceptions. No prior work has brought FiLM conditioning into per-flow online routing, leaving differentiated topology perception and policy learning under mixed traffic an open problem.

To close this gap, we propose QoS-FiLM. The core idea is to inject the QoS intent vector of each flow into every message-passing layer of a GNN via FiLM, so that the same physical topology is encoded into differentiated node embeddings under different intents. We then formalize online mixed-service routing as a Markov decision process (MDP) and train a non-myopic routing policy with proximal policy optimization (PPO). The main contributions of this paper are as follows:
\begin{itemize}
    \item \textbf{A flow-aware graph reinforcement learning architecture for mixed-service online routing.}
    The architecture couples a flow-conditioned graph encoder, a path-aware scorer, and an independent value network, all trained end-to-end via proximal policy optimization. It resolves the core contradiction that the same topology must yield differentiated routing decisions for flows with different QoS intents.

    \item \textbf{A per-flow QoS-intent-driven FiLM conditioning mechanism and network design.}
    The per-flow intent vector generates affine modulation parameters for each GNN layer, rendering the message-passing process flow-aware. We further design complementary modules, including LSTM-based path embedding pooling, intent-path joint scoring, and intent-conditioned attention graph pooling, which together realize a complete flow-conditioned decision chain from topology encoding to path selection.

    \item \textbf{Dual-space interpretability validation.}
    Experiments show that the proposed architecture significantly outperforms classical baselines in cumulative utility, bandwidth satisfaction ratio, and end-to-end latency. A dual-space analysis in the representation space (hidden feature separability) and the decision space (path selection differentiation) further confirms that merely changing the QoS intent triggers differentiated routing decisions on the same topology and the same source-destination node pair, validating the necessity of FiLM conditioning at the mechanism level.
\end{itemize}

The remainder of this paper is organized as follows. Section~II reviews related work on online routing and conditioned graph representation learning. Section~III formalizes online mixed-service routing as a Markov decision process and defines the optimization objective. Section~IV details the QoS-FiLM architecture, including the flow-conditioned graph encoder, the path-aware scorer, the independent value network, and the per-layer FiLM conditioning mechanism. Section~V reports the experimental design, main results, ablation analysis, and interpretability validation on the GEANT, NOBEL-EU, and ABILENE backbone topologies. Section~VI concludes the paper and outlines future work.
